In this appendix we further illustrate the possible use of an alternate table joined with the {\tt ivoa.obscore} table as a method for accessing {\bf response-function}s and data products that require queryable attributes that are not present in the {\tt ivoa.obscore} table,  as discussed in \S~\ref{sec:respaccess}. Further scientific input from \gls{HEIG} domain experts, developed from additional use cases that focus on exploring the appropriate set of attributes needed to support such queries, will be needed to further explore and generalize this solution.

In order to handle the various possible cardinality relationships between {\bf response-function} and {\bf hea-event-list}  datasets, foreign keys must be defined in the {\tt ivoa.response} tables that will allow {\tt  JOIN} operations between those tables and the {\tt ivoa.obscore} table.

The {\tt ivoa.response} table uses a specific resp\_dataproduct\_type column which is compliant to the response type vocabulary : \url{https://www.ivoa.net/rdf/response-type/} and has its own resp\_obs\_publisher\_did column as well, in order to distinguish which response matches  which ObsCore data set in the {\tt ivoa.response} table. 

\subsection*{Proposal for a Response Table {\tt ivoa.response}}

{\small
 \begin{longtable}{|m{0.25\linewidth}|m{0.1\linewidth}|m{0.10\linewidth}|m{0.41\linewidth}|}
\hline
{\bf Column Name} & {\bf Unit} & {\bf Type} & {\bf Description}\\\hline
%obs\_collection & unitless & String & Name of the data collection \\\hline
% to bind a response with an observation data set
{\em obs\_id} & unitless & String & Related observation ID (foreign key) \\\hline
{\em obs\_publisher\_did} & unitless & String & Related dataset ID (foreign key) \\\hline
% copied from related observation if needed
{\em target\_name} & unitless & String & Astronomical object observed, if any\\\hline
{\em facility\_name} & unitless & String & Facility name used for observed data products \\\hline
{\em instrument\_name} & unitless & String & Instrument name used for observed data products \\\hline
% response specific fields 
{\em resp\_product\_type} & unitless & String &  Product type as defined in the response-type vocabulary\\\hline
{\em resp\_publisher\_did} & unitless & String & Dataset identifier given by the publisher\\\hline
{\em resp\_access\_url} & unitless & String & URL used to access (download) the response dataset\\\hline
{\em resp\_access\_format} & unitless & String & File content format (see in Obs\-Core MIME Types)\\\hline
% response spatial coverage
{\em resp\_s\_ra} & deg & double &  ICRS reference right ascension, for the region where the response-product applies \\\hline
{\em resp\_s\_dec} & deg & double & ICRS reference declination, for the region where the response-product applies\\\hline
{\em resp\_s\_region} & unitless & String & Region of interest for response use (expressed in ICRS frame)\\\hline
{\em resp\_t\_intervals} & unitless& String (TMOC) & Time coverage for response use \\\hline
%response energy coverage
{\em resp\_energy\_min} & eV & double & Energy band minimal value for response use  \\\hline
{\em resp\_energy\_max} & eV & double & Energy band maximal value for response use  \\\hline

\caption{Proposal for a standard response table :  with appropriate further study the columns identified herein could form the basis for a recommended base set of columns for the  {\tt ivoa.response} table.}
\label{tab:response_table}
\end{longtable}
}
%\end{center}
%\end{table}

\subsection*{Example queries}
\noindent Implementing an {\tt ivoa.response} table similar to this would allow queries such as the following: \\
\subsubsection*{Retrieve psf response for {\em obs\_id} and position criterium }
\noindent Find all datasets satisfying:
\begin{enumerate}[(i)]
 \item Position inside 3 arcmin from (83.6324, $+22.0174$),
 \item Response product\_type = ``psf'',
 \item obs\_id = ``1527'',
 \item obs\_collection = ``HESS''.
\end{enumerate}

When using the response table defined in Table \ref{tab:response_table}, we can join the {\tt ivoa.obscore} table with the {\tt ivoa.response} table on 
the two keys  {\em obs\_id\/} and {\em obs\_publisher\_did\/} for instance, as in: 
{\small
\begin{verbatim}
SELECT o.obs_publisher_did, o.obs_id, s_ra, s_dec, s_fov, t_min, t_max, 
energy_min, energy_max, access_url, access_format, 
resp_publisher_did, r.obs_publisher_did, r.obs_id, resp_access_url, resp_access_format, 
resp_energy_max, resp_energy_min
FROM ivoa.obscore as o
NATURAL JOIN  ivoa.response as r
WHERE
(target_name = ’Crab’ OR target_name = ’M1’ OR
CONTAINS(POINT(s_ra, s_dec), CIRCLE, 83.6324, +22.0174, 0.083333) = 1)
AND (resp_product_type = 'psf')
AND (o.obs_id = '1527')
AND (o.obs_collection = 'HESS')
\end{verbatim}
}
If additional constraints are required, for instance on time, then we can add another set of columns  in {\tt ivoa.response}({\em e.g.\/}) {\em resp\_t\_min\/}, and {\em resp\_t\_max\/} and constrain the time interval by adding a clause like:
{\small
\begin{verbatim}
AND (resp_t_min > 56000.0 and resp_t_max < 56001.5) 
\end{verbatim}
}

\subsubsection*{Retrieve edisp response for data sets constrained by {\em obs\_id} or position and by {\em scan\_mode}}
{\small
\begin{verbatim}
SELECT o.obs_publisher_did, o.obs_id, s_ra, s_dec, s_fov, t_min, t_max, 
energy_min, energy_max, access_url, access_format, scan_mode,
resp_publisher_did, resp.obs_id, resp_access_url, resp_access_format, 
resp_energy_max, resp_energy_min, 
FROM ( SELECT o.obs_publisher_did, o.obs_id, scan_mode 
     FROM ivoa.obscore 
     NATURAL JOIN  ivoa.obscore-hea  
     WHERE
     (o.obs_collection = 'IACT') 
     AND (scan_mode=raster-map)
     AND (o.obs_id = '1976')
     OR CONTAINS(POINT(s_ra, s_dec), CIRCLE, 83.6324, +22.0174, 0.083333) = 1) as s
JOIN ivoa.response ON 
resp.obs_id =s.obs_id
WHERE (resp_product_type = 'edisp')
% s is the list of observations satisfying position or obs_id and scan mode criteria
\end{verbatim}
}% end \small



